\chapter[Who Owns and Displays Art?]{Who Owns Art, and Who May Display It?}
\section{A Ticket That Cost Nothing}
First, pick up an object. Not a bronze or a famous painting, but a piece of paper---more precisely, a ticket you cannot actually obtain.

London's British Museum charges no admission. Walk in, pass security, and you can stand before the Rosetta Stone. Inscribed with three scripts, it unlocked ancient Egyptian hieroglyphs; archaeologists around the world recognize it. Looking costs nothing. Farther inside, the Parthenon sculptures, Egyptian mummies, and Chinese Tang sancai pottery are all free too. There can scarcely be a better bargain: some of humanity's most precious objects from thousands of years, and you can look as long as you please without spending a pound.

But pause and examine that word ``free.'' Something free to you is not necessarily without a price; someone else has paid the bill. The museum's land, lighting, temperature and humidity control, guards, and conservators cost astronomical sums each year. A larger bill lies deeper: how did these objects get here? When they left their homelands, was a price paid? To whom? Or was nothing paid at all?

So we begin with this nonexistent free ticket. Behind it lies a question: who paid for everything we view without charge? And who decides what ownership and display mean?

\section{An Afternoon in Gallery 33}
Now enter a scene.

Time: a July afternoon. Rain has just passed over London, leaving the pavement wet. Place: Gallery 33, the British Museum's Chinese gallery. It is a long room, dimly lit to protect light-sensitive silk and paper. The air smells of carpet and old building. Audio guides murmur, footsteps pass, and shutters click softly between glass cases.

A study-tour group of Beijing secondary-school students crowds in. The last girl, fifteen-year-old Xiao Yu, grips a gallery map softened by her sweat.

First she sees a life-size Liao-dynasty sancai luohan, or Buddhist arhat, seated cross-legged with slightly knitted brows, as though considering a problem unresolved for a thousand years. Its label says Yixian, Hebei, China; acquired in the early twentieth century. Farther inside, a case holds a reproduction of a long scroll: a Tang copy of Gu Kaizhi's Admonitions of the Instructress to the Court Ladies. The actual painting, vulnerable to light, is displayed only a few weeks each year. Their teacher whispers, ``The real painting is in this museum. It is one of our national treasures.''

Xiao Yu stares at ``acquired'' for several seconds. What a strange word. It tells no lie, but says almost nothing. How did the luohan travel from a cave in Yixian to London? Who moved it, sold it, agreed to it? That little line of print resembles a tightly shut door concealing an entire untold route.

She looks up and asks, ``Teacher, can they still go home?''

The teacher does not answer immediately. Nearby, a local British primary-school child leans against the same case while a teacher points out the luohan, drawing quiet exclamations from the children. Suddenly Xiao Yu notices something uncomfortable. These national treasures, so far from home, are well cared for and widely viewed. Their lighting is more professional than that of many objects back in China. They seem to lead respectable lives abroad---but not lives they chose.

Remember this gallery and the word ``acquired.'' Most of the chapter begins behind the door that word conceals.

\section{What Gives It the Right to Be Here?}
First lay out the question without rushing to answer.

Xiao Yu's question sounds simple: should these things be returned? Dig one layer deeper, however, and it becomes four overlapping questions.

First, how did the object arrive? This is a factual question answered through evidence: invoices, accounts, letters, legal documents, and participants' recollections.

Second, who owns it now? This is a legal question. Under the museum's national law, the answer is often clear: the museum owns it, as the property records explicitly state.

Third, where does it belong? This question extends far beyond law into cultural belonging. Does a luohan that sat in a Yixian cave for eight hundred years have a connection with Hebei's land and the descendants of its people that law cannot govern?

Fourth, who may interpret it? Even if you accept its presence in London, who writes its label and tells its story? Is it presented as a masterpiece of human civilization or evidence of colonial plunder? The same luohan, explained in two ways, sends visitors away with entirely different worlds.

People often quarrel because they collapse these four questions into one. Calling something humanity's shared heritage addresses the third layer. Replying that it was illegally exported addresses the first and second. Claiming superior preservation states a fact but tries to answer what ought to happen, substituting one question for another.

We will separate these layers. Once we do, ownership and display will prove far larger issues than you expect. They concern not only museums but conservators' tools, painters' borrowings, the memes on your phone, and the enormous argument now unfolding about artificial intelligence.

Before continuing, choose a position: if removing an artifact followed procedures unobjectionable under the laws of its time, should it nevertheless be returned today?

\section{Defending Museums and Demanding a Return}
Let us present both sides fully---not select their most ridiculous arguments to mock.

\textbf{Position one: defending the universal museum.}

Make the British Museum's case as strongly as possible. Chinese readers often consider this side the loser, and losers' arguments frequently go unheard.

First, rescue and preservation. Consider the Parthenon sculptures. When Elgin, Britain's ambassador to the Ottoman Empire, removed them in the early nineteenth century, Athens was under Ottoman rule. The temple had suffered warfare and a devastating explosion during the preceding century or more; sculptures left there would continue to weather or be stolen and sold. Similar reasoning applies to many Chinese artifacts: amid early twentieth-century warfare, numerous objects survived precisely because they left.

Second, shared heritage. These objects belong to humanity, not merely one modern nation. Gathering Greek sculpture, an Egyptian stone, and a Chinese luohan under one roof lets anyone buy a plane ticket and survey human civilization in a day---an educational achievement. Free admission also lets rich and poor encounter the same originals.

Third, legality. Elgin claimed a permit from the Ottoman authorities then legally ruling Greece. Many Chinese artifacts were bought on the market in transactions not then unlawful. Applying today's laws to acts two hundred years ago is retroactivity, itself prohibited under modern rule of law.

Fourth, custodial capacity. Some objects left in their places of origin might long ago have succumbed to war, theft, or poor preservation. This is not merely hypothetical; many instances exist.

These arguments are substantial, and many sincere scholars hold them.

\textbf{Position two: let them go home.}

Now hear the other side just as fully.

First, consent must be examined within the power structures of its time, including those now obsolete. The Ottoman permit came from occupiers, not Greeks. If a landlord sold a tenant's house, would the tenant's signature count? As for lawful purchase, consider the starkest example: the Benin Bronzes. In 1897, British troops attacked the Kingdom of Benin, in present-day Nigeria, in the name of retaliation, burned its palace, and carried to Europe more than a thousand bronzes and ivory carvings accumulated by its royal court over centuries. They were dispersed among dozens of museums. This was removal after military victory, not purchase. Procedurally, every step occurred under colonial-war rules---written by those holding the guns.

Second, contextual integrity. Parthenon sculptures are not a framed painting; they are architectural elements sawn from a temple's body. Tear a book in half, leave one half in Athens and the other in London, and neither place can read more than a fragment. The Acropolis Museum opened in Athens in 2009, at the site's foot, with spaces prepared for the sculptures. The claim that there is nowhere to put them has expired.

Third, custodial-capacity arguments are circular. They assume source communities can never protect their own possessions, although that incapacity was often produced by the same colonial order: first disrupt a society, then remove its objects because it supposedly cannot manage them.

Fourth comes an easily overlooked voice: not everything exists to be viewed. Some Indigenous North American communities regard certain sacred objects as unsuitable for display or photography. They belong to ritual, not galleries. For these communities, the right to display matters as much as ownership, sometimes more painfully. Offering an object for worldwide museum appreciation is therefore not neutral. It assumes a particular value: an object's worth lies in being seen.

The restitution list stretches across continents. For more than a century, Egypt has sought Nefertiti's painted bust from Germany and the Rosetta Stone from Britain. Greece has pursued the Parthenon sculptures for two hundred years. Residents of Chile's Easter Island seek an ancestral stone figure. Ethiopia seeks royal treasures taken by British troops from Maqdala. This is not a grievance held by one or two countries, but an old account as large as a world map, because the original removals were a world-map-sized undertaking.

\textbf{Recent developments deserve attention}: actual returns have begun. In 2022, Germany transferred ownership of hundreds of Benin Bronzes to Nigeria, and some American and British institutions have returned portions of their collections. The dispute is no longer merely verbal. Yet British law restricts the British Museum's ability to dispose of collections unless the law changes. The final gate returns us to who wrote the rules.

Both sides claim to protect human civilization. They differ over what that phrase means: objects behind glass for everyone to view, or living things rooted in land, communities, and memory.

\section[Thinking Bridge: Tracing Provenance]{Thinking Bridge: Trace an Artifact's Provenance}
Emotion and allegiance cannot settle restitution debates. Museum studies and law offer a portable tool: \textbf{provenance research}, tracing a complete chain of an object's origins and ownership, its life written like a resume.

The chain has roughly five stages:

\textbf{Creation $\rightarrow$ Placement or use $\rightarrow$ Departure from home $\rightarrow$ Transfers $\rightarrow$ Acquisition and display}

Its power comes from asking four questions at every link:

\begin{enumerate}
\item How did it move: voluntary gift, market transaction, war booty, smuggling, or unclear circumstances?
\item Who consented? Was that person or institution entitled to decide for everyone?
\item What were the power relations? Were the parties equal? Were guns nearby?
\item Who benefits now? Where do admission revenue, prestige, research rights, and tourism income flow?
\end{enumerate}
Try this with Xiao Yu's Liao sancai luohan. Removed from a Yixian cave in the early twentieth century, it passed through dealers and was sold overseas. How did it move? Market purchase. Who consented? Local intermediaries, but could they dispose of a community's shared heritage? Moreover, there were originally several luohans, and some suffered damage during illicit removal. Power relations? This was an era of warlord conflict, national weakness, and poverty, with sellers and buyers on entirely different scales. Who benefits? Today's museum gains visitors and prestige; the place of origin gains a name on a label. After these questions, ``lawful purchase'' looks much thinner. Every link may have a receipt while the whole chain remains unjust.

The tool also prevents misjudgment. Consider the Shosoin repository at Todaiji in Nara, Japan. It holds exquisite Tang Chinese objects: lutes, bronze mirrors, aromatics, and documents. Your first thought might be more stolen treasures, but tracing provenance reveals that most arrived through Japanese missions to Tang China, trade, and gifts. Japan's imperial household recorded and preserved them as heirlooms for over a thousand years. Both cases involve Chinese artifacts abroad, yet one chain bears blood and fire, another seals and gift lists. \textbf{``Every Chinese artifact overseas was stolen'' is an easy judgment that fails a provenance test.} The tool's value is precisely that it forbids one blanket sentence about a whole cabinet.

Provenance can leave the museum too. An online image claiming to be a genuine Qing photograph, an ancient original sound recording, an ancestral secret recipe: ask where it came from, whose hands it passed through, and who benefits. This is self-defense for both the museum age and the digital age.

\section{Two Bandits in a Letter}
Now read a brief primary source. In 1861, French writer Victor Hugo received a letter from a French officer who had participated in the invasion of China, asking his view of the expedition. His reply became known as the Letter to Captain Butler. It says, in paraphrase with its key sentences in their familiar translation:

\begin{quote}
In a corner of the world stood a wonder of humanity: the Summer Palace. One day, two bandits from Europe entered it. One looted, the other set fires. The two bandits whom history will judge were named France and England.
\end{quote}
The letter matters not simply for its fierceness but for its author. Hugo was French, among the winning side's most celebrated writers. If victors write history, he should have praised the expedition. Instead, he uncompromisingly called his own country's conduct robbery: in paraphrase, ``We Europeans are civilized; the Chinese are barbarians in our eyes. This is what civilization did to barbarism.''

The letter teaches us to distinguish three things. What happened: in 1860, Anglo-French forces looted and burned the Old Summer Palace during the Second Opium War, dispersing many artifacts overseas. That is evidential fact. How contemporaries understood it: Hugo shows that contemporaries never spoke with one voice. Even on the winning side, some refused to call plunder victory. What people then said cannot stand as a unanimous declaration for either side. How we evaluate it: accounts in Chinese and French textbooks, museums, and diplomacy continue to change. This history is still being written.

One small but crucial caution: the letter's fame does not mean it has made our judgment for us. It demonstrates that someone recognized the deed's character more than 160 years ago. It does not answer what should happen to the artifacts today, nor can Hugo answer for people living after his time. Evidence can narrow a dispute, but cannot make your value choices.

Incidentally, the objects' later paths form a living provenance textbook. Some reached France and are now displayed in Fontainebleau's Chinese Museum. Others passed through British and European auction houses into public and private collections. Some remain untraced. In recent years, purchasers have donated animal heads and other Old Summer Palace objects back to China. At the end of the same chain, a new link is appearing: homecoming.

\section{Restoration: Rescue or Cosmetic Surgery?}
Change battlefields. After ownership comes a quieter conflict, equally lacking standard answers: the conservator's workbench.

Start with undisputed facts: an excavated bronze lies in dozens of pieces, an old painting's pigments are flaking away, insects have hollowed a wooden pagoda's beams. Do nothing and they will die. Conservation is necessary.

Disagreement begins over how far to restore. Again, the same facts receive different interpretations. Set out the two sides.

\textbf{Interpretation one: preserve evidence through minimal intervention.}

This side treats an artifact first as evidence. Every crack, repair, and later paint layer carries historical information, like tree rings. The conservator's sole task is to stop the bleeding, preventing further deterioration, never making it whole again. Leave a Buddha's missing ear missing. Where additions are unavoidable, distinguish them from the original. The Chinese principle of restoring the old as old broadly expresses this: afterward, it should remain the old object, not become a beautiful new one.

This position is honest: future generations can distinguish history from our craftsmanship. Its cost is that visitors often encounter incomplete works and must fill the gaps imaginatively.

\textbf{Interpretation two: restore legibility so the work can be understood.}

The other side treats the artifact first as art. Art was made for people. If damage makes its original greatness wholly unreadable, preservation saves only mute fragments. The most famous battlefield is Michelangelo's Sistine Chapel ceiling. In the 1980s and 1990s, the Vatican undertook extensive cleaning and restoration. Removing five hundred years of grime and old coatings revealed startlingly vivid colors. Some experts celebrated the master's rebirth; others lamented overcleaning that, they believed, removed final shadow layers Michelangelo himself had applied. They said they now saw the conservators, not Michelangelo. Both sides faced the same ceiling.

\textbf{Now consider an example that shakes the entire debate.} Ise Jingu in Japan's Mie Prefecture, Shinto's most sacred shrine, began a practice over thirteen hundred years ago: every twenty years, shikinen sengu rebuilds the entire sanctuary in identical form with new timber on an adjacent site, then dismantles the old one. Through a European stone-monument lens, this seems disastrous: the original is destroyed every twenty years, and no timber in the current sanctuary is older than twenty. Yet the Japanese see no puzzle. The true original is not that batch of wood but the form, craft, and people transmitting the craft. Rebuilding every twenty years prevents that skill from dying: experienced masters and apprentices build it again from the ground up, keeping the technique alive.

This contrast is unsettling. Preserving the original seems common sense, but what counts as original is itself culturally answered. In a stone-centered tradition, it is the material; in a timber tradition, it may be the form. Neither is wrong, but neither is the uniquely correct custodian.

Together these positions suggest a modest conclusion: restoration is half science, half standpoint. Science concerns materials, chemistry, and techniques; standpoint concerns the past we wish to leave future generations. Good conservators do not pretend the latter half is absent. They recognize every cut as a choice and record their choices for later inspection.

\section[What Remains of the Original?]{Deeper Challenge: What Remains of the Original?}
Now for the chapter's theoretical heavyweight. In 1935, German thinker Walter Benjamin wrote The Work of Art in the Age of Mechanical Reproduction. Short though it is, its influence reaches every art museum today. He proposed the concept of \textbf{aura}, also rendered as a halo or spiritual radiance.

In paraphrase, an original artwork has a unique here-and-now existence. It is not simply painted cloth, but that particular cloth: Leonardo touched it, it stayed in Florence, was stolen, and hangs on a particular Louvre wall. Five centuries of viewing, stories, and wear have accumulated within it. That presence brewed from time, place, and history is aura. Photography and film made reproduction easy; anyone could take home the Mona Lisa on a postcard. As copies proliferate, aura begins to fade, and art's social function changes. A mysterious object emerging from worship and ritual becomes something exhibited, consumed, and discussed.

This theory explains an everyday oddity. Your phone's Mona Lisa image has higher resolution than the view beyond three rows of people in the Louvre. Yet millions still fly to Paris each year, queue for hours, and stand before it for thirty seconds. They are not seeking an image, available freely. They are making a \textbf{pilgrimage} to the aura reproduction cannot carry away. Benjamin foresaw aura's decline, but the Louvre's queues suggest it is more resilient than expected.

The theory also illuminates restoration. Conservators claim to protect aura by preventing the original's death. Critics call excessive restoration \textbf{counterfeiting aura}: giving an old object cosmetic surgery and inviting continued reverence for its five-hundred-year age. Ise Jingu offers a third answer: aura need not inhabit material; it may live in craft and ritual, relit every twenty years. These practices offer three answers to where aura lives.

The theory has limits, and the opposing case deserves its full force. Benjamin's attitude toward aura's decline was complex, even partly welcoming: bringing art down from its pedestal lets small-town children see the Sistine ceiling in a picture book. Reproduction democratizes as well as dilutes. Reproductions in county-town classrooms exemplify something the age of aura could not achieve. Neither absolute supremacy of originals nor unqualified celebration of copies is Benjamin's conclusion. His real warning is: \textbf{when reproduction changes our relationship to art, all our inherited ideas of ownership, authenticity, and value require reexamination.}

We now stand amid a reproductive revolution ten thousand times fiercer than photography. Digital files copy without cost or degradation; a first and hundred-millionth copy are identical, nearly undoing the concept of an original. AI goes further: after seeing hundreds of millions of images, a model can produce a new image closely resembling a painter's style. How do we discuss originality then? Benjamin did not live to see this, but his question lands squarely on us.

Another development deserves attention: digital restoration. Technology can generate how a damaged statue might once have looked or infer a faded mural's former colors. This is essentially the conservator's old choice, changed from physical additions to algorithmic guesses. Its benefit is leaving the object untouched; its danger is that the guessed original may resemble contemporary tastes more than we imagine. Preservation or a new version: the old examination has changed rooms, not ended.

\section{Today: Every Image You Post Enters a Dispute}
This centuries-old problem sits in your pocket. First consider two sets of ongoing developments, then apply the chapter's tools.

\textbf{First: copyright, the public domain, and re-creation.}

Modern copyright began with Britain's 1710 Statute of Anne. It did something novel: assigned authors rights over their works for a limited term. Afterward, works entered the \textbf{public domain}, available for anyone to print, adapt, or recreate without payment or permission. At its founding, the United States placed the same idea in its Constitution. Article I, Section 8 empowers Congress:

\begin{quote}
To promote the Progress of Science and useful Arts, by securing for limited Times to Authors and Inventors the exclusive Right to their respective Writings and Discoveries.
\end{quote}
Focus on limited times and promoting progress. From the outset, copyright was a bargain, not a natural ownership right: society gave authors temporary exclusivity to earn a living, in exchange for eventual return of the work to everyone. Anyone may adapt or perform Du Fu, Shakespeare, or Beethoven, not because law cannot reach them, but because their works long ago became public property shared by humanity. Disney's Mickey Mouse, first appearing in 1928's Steamboat Willie, entered the US public domain on January 1, 2024: that earliest Mickey can now be legally adapted.

Now inspect your phone. You turn anime screenshots into memes, write new lyrics for a favorite song and post them, or draw fan art using game characters. Each has legal complexities, but three ethical questions help: is the creator acknowledged? Are you making money? Does your re-creation promote the original or substitute for it in the market? Most casual derivative creations are safe under these questions, but everyone-does-it is never a reason; the three questions are.

The newest problem comes from AI. Large models train on hundreds of millions of online artworks, photographs, and articles, with most creators uninformed, unasked, and unpaid. The established fact is that artists, writers, and media organizations have brought multiple copyright cases in courts around the world, with law not yet settled. Viewed through provenance, does this resemble the museum problem enlarged a hundred million times? Works are taken; those taking them invoke humanity's progress; those losing them have no chance to consent or refuse. History does not repeat, but it rhymes.

Zoom closer, into school. At a New Year show, a class rents Miao silver-adorned ceremonial outfits online and wears them for street dance. The audience applauds, but an online comment asks whether silver ornaments associated with solemn Miao occasions should become performance props. Another replies that beauty, liking, and absence of ill intent constitute cultural exchange. Even a classroom stage cannot escape the three questions. Nor can illustrations grabbed online for a group display: were they credited? Cropping a classmate's watermark from a photograph and posting it to your own public account differs from removing a mask from Benin's palace by a hundred orders of magnitude, yet violates the same small rule: do not erase origins.

\textbf{Second: cultural appropriation and cultural exchange.}

In 2022, Dior released a skirt whose silhouette many Chinese consumers identified as closely resembling the traditional mamianqun, or horse-face skirt. The brand described it as its own signature silhouette without mentioning Chinese origins, provoking substantial protest. In 2015, Boston's Museum of Fine Arts held Kimono Wednesdays, inviting visitors to try on kimono and pose beside a Monet painting. Fierce disagreement followed: some called it consumption of another's culture, while some Japanese people in America expressed no objection. Even members of the same community disagreed.

These cases show that cultural appropriation has no universal alarm, but a usable framework returns to three questions: \textbf{are origins acknowledged, can the people concerned say no, and who receives the benefits?} The Dior incident hurt not because it used Chinese elements, but because it erased origins and claimed others' creativity as its own signature. Conversely, condemning every borrowing as appropriation also harms. Picasso drew inspiration from African masks and opened a new era in modern art; jazz, pizza, and jeans all resulted from cultural movement. Japanese custodians have treasured Tang lutes at the Shosoin for thirteen hundred years while openly recording their origins: that is exchange. If someone instead takes the same lute's design, erases its origins, and charges patent fees to its source community, something else has occurred. The distinction concerns not borrowing itself, but equal power and acknowledgment of the account.

\section{For You: A Museum Half Empty}
Return to Xiao Yu's gallery.

Imagine an international convention taking effect tomorrow. Every artifact acquired through plunder or deceit, every provenance chain marked by blood and guns, requires unconditional return. Half the British Museum's galleries empty, a corner of the Louvre empties, and entire walls on Berlin's Museum Island stand bare. Meanwhile, crowds fill museums in Cairo, Athens, Beijing, and Benin City. For the first time, children see objects from ancestral stories under their own cities' lights.

Is this a better world?

You have arguments for yes: justice is fulfilled, torn contexts repaired, and the authorship of rules finally changes. Nigerian and Egyptian children who could never travel abroad can encounter their own greatness at home.

You have arguments for worse too: globally dispersed collections provide a distinctive education, and a small-town young person's chance to survey civilization in one London day disappears. Some source locations genuinely lack adequate preservation conditions. And where does return end? How far back should provenance reach: a hundred years, five hundred, the Crusades? Might it become an account that can never be settled but makes every museum anxious?

A sharper question hides behind both answers. If art can belong to all humanity and to a particular people, with both claims sounding just, \textbf{might ownership itself be too small a concept to contain art?}

What is your answer? Give reasons. And next time you enter a museum, try this: before looking at the object, read the tiny provenance line on its label and continue the chain it leaves unfinished. That is the ticket this book truly wants to give you.
